\section{Research Gap and Novelty Positioning}
\label{sec:research_gap}

A systematic analysis of our literature corpus (encompassing 34 foundational and contemporary studies across cognitive architectures, external memory networks, agent memory stores, and context management frameworks) reveals five critical, unaddressed gaps that limit the reliability and auditability of long-running autonomous agents. Among the 34 papers reviewed in our literature corpus, only 5 (14.7\%) explore any versioning concepts, and none integrate automated DAG commit tracking with budget-constrained knapsack context compilation.

\subsection{Gap 1: Destructive Mutation vs. Immutable Versioning}
In mainstream persistent agent memory architectures (e.g., MemGPT~\cite{packer2023memgpt}, Mem0~\cite{deshmukh2024mem0}), memory updates are executed as destructive in-place text replacements or SQL \texttt{UPDATE}/\texttt{DELETE} operations. Once modified, prior states are obliterated. If an agent misbehaves or suffers a prompt injection attack, developers cannot reconstruct the agent's exact epistemic state at that moment.

\textit{Counterexample Analysis:} Git Context Controller (GCC)~\cite{wu2025gcc} and GitOfThoughts~\cite{shekar2025gitofthoughts} use Git commits, but require the LLM itself to manually invoke tool commands to write and commit notes. This imposes significant token overhead and relies on the model remembering to manage its own version control. In contrast, ManiGit provides an \textbf{automated, transparent infrastructure layer}: incoming messages are automatically parsed, classified, and committed to an immutable commit DAG in SQLite without requiring agent tool calls or prompt tokens.

\subsection{Gap 2: Contradiction Handling and Ghost Context}
In vector-based memory systems, text chunks are indexed as timeless semantic points. When user constraints change across sessions (e.g., migrating from PostgreSQL to MongoDB), vector similarity search retrieves both the obsolete initial choice and the new requirement. This introduces \textit{ghost context} into the prompt, forcing the language model to arbitrate conflicting premises. Recent evaluations on STALE~\cite{stale2025benchmark} and LongMemEval~\cite{wu2024longmemeval} confirm that frontier models fail on implicit contradictory premises in over 40\% of trials.

\textit{Partial Solutions:} Generative Agents~\cite{park2023generative} applies recency weighting, which merely dampens older memories without removing them. MemoryBank~\cite{zhong2023memorybank} applies exponential decay based on time rather than logical contradiction. In contrast, ManiGit implements \textbf{explicit contradiction detection and supersession}. When a new fact is ingested, active memories are evaluated for semantic conflict. Superseded facts are automatically transitioned to the \texttt{ARCHIVED} tier, removing them from active prompt compilation while preserving their provenance in the commit graph.

\subsection{Gap 3: Point-in-Time Historical Reconstruction}
When auditing long-running agents or answering retrospective questions (e.g., ``What was our deployment target prior to commit $C_k$?''), existing persistent systems evaluated in our literature review generally do not provide the combination of versioned memory state and point-in-time reconstruction, as past states have either been overwritten or blurred into an uncurated flat log.

\textit{Counterexample Analysis:} While Git provides checkout semantics for source code, prior agent architectures do not provide point-in-time state reconstruction for conversational context memory under budget constraints. ManiGit provides \textbf{first-class historical checkout}: given any historical commit hash $C_k$, ManiGit reconstructs the exact active memory set as it existed at $C_k$ and compiles a budget-compliant context window for that exact historical moment, achieving 100\% recovery accuracy across the evaluated benchmark checkout scenarios.

\subsection{Gap 4: Token-Budget Context Compilation (Memory $\ne$ Context)}
Existing systems frequently conflate persistent memory storage with prompt context assembly. Naive in-context approaches stuff as many retrieved chunks as possible until prompt limits are reached. However, empirical studies confirm that long, bloated contexts degrade model attention~\cite{liu2024lost,dong2023bamboo}.

\textit{Counterexample Analysis:} Classic information retrieval has formulated sentence extraction as a knapsack problem for multi-document summarization, but has not integrated it with dynamic, versioned memory tiers. ManiGit formalizes context construction as a \textbf{bounded knapsack optimization problem} subject to a hard token budget $B$:
\begin{equation}
\max_{C_t \subseteq \mathcal{M}_t} \sum_{m_i \in C_t} U(m_i, q_t) \quad \text{subject to} \quad \sum_{m_i \in C_t} L(m_i) \le B
\end{equation}
ManiGit addresses this via a deterministic greedy score-density compiler that balances query relevance, composite priority, and memory tier weighting, maintaining zero budget violations ($\sum_{m_i \in C_t} L(m_i) \le B$) across all evaluated trials.

\subsection{Gap 5: Provenance and Auditable Why Reports}
When an agent acts on a stored memory, standard vector stores offer no visibility into why the fact was retained, which user message generated it, or what confidence score was assigned.

\textit{ManiGit Positioning:} Every memory fact in ManiGit maintains bidirectional pointers to the originating raw message ID and turn index, accompanied by a structured \textbf{Why Report} documenting the classification rule, LLM judge confidence score, and supersession target.
